\documentclass[11pt,a4paper]{article}

% ── Codificación y fuentes ────────────────────────────────────────────────────
\usepackage[utf8]{inputenc}
\usepackage[T1]{fontenc}
\usepackage{lmodern}
\usepackage[spanish]{babel}

% ── Geometría y espaciado ─────────────────────────────────────────────────────
\usepackage[top=2.5cm, bottom=2.5cm, left=2.8cm, right=2.8cm]{geometry}
\usepackage{setspace}
\setstretch{1.15}
\usepackage{parskip}

% ── Tablas ────────────────────────────────────────────────────────────────────
\usepackage{booktabs}
\usepackage{tabularx}
\usepackage{array}
\usepackage{longtable}
\usepackage{enumitem}

% ── Colores y cajas ───────────────────────────────────────────────────────────
\usepackage[dvipsnames]{xcolor}
\usepackage{tcolorbox}
\tcbuselibrary{skins, breakable}

% ── Cabeceras y pies de página ────────────────────────────────────────────────
\usepackage{fancyhdr}
\usepackage{lastpage}
\pagestyle{fancy}
\fancyhf{}
\fancyhead[L]{\small\textbf{Informe de Evaluación} --- Física para Bioanalistas}
\fancyhead[R]{\small krischel bravo 32449538}
\fancyfoot[C]{\small Página \thepage\ de \pageref{LastPage}}
\renewcommand{\headrulewidth}{0.4pt}

% ── Hiperenlaces ──────────────────────────────────────────────────────────────
\usepackage[colorlinks=true, linkcolor=NavyBlue, urlcolor=NavyBlue]{hyperref}

% ── Paleta de colores personalizados ─────────────────────────────────────────
\definecolor{desempeno_excelente}{HTML}{1A6B2F}
\definecolor{desempeno_bueno}{HTML}{2E5A9C}
\definecolor{desempeno_suficiente}{HTML}{8B6914}
\definecolor{desempeno_deficiente}{HTML}{C0392B}
\definecolor{desempeno_insuficiente}{HTML}{7B241C}
\definecolor{grisFondo}{HTML}{F5F5F5}
\definecolor{azulTitulo}{HTML}{1B2A6B}
\definecolor{verdeFortaleza}{HTML}{1A6B2F}
\definecolor{naranjaMejora}{HTML}{A04000}

% ── Estilos de cajas ─────────────────────────────────────────────────────────
\tcbset{
  cajaTitulo/.style={
    enhanced, breakable,
    colback=azulTitulo!8, colframe=azulTitulo,
    fonttitle=\bfseries\large, coltitle=white,
    attach boxed title to top left={yshift=-2mm, xshift=4mm},
    boxed title style={colback=azulTitulo},
    top=6pt, bottom=6pt, left=8pt, right=8pt,
  },
  cajaFortaleza/.style={
    enhanced, breakable,
    colback=verdeFortaleza!6, colframe=verdeFortaleza!60,
    fonttitle=\bfseries, coltitle=verdeFortaleza,
    top=4pt, bottom=4pt, left=8pt, right=8pt,
  },
  cajaMejora/.style={
    enhanced, breakable,
    colback=naranjaMejora!6, colframe=naranjaMejora!60,
    fonttitle=\bfseries, coltitle=naranjaMejora,
    top=4pt, bottom=4pt, left=8pt, right=8pt,
  },
  cajaRecomendacion/.style={
    enhanced, breakable,
    colback=grisFondo, colframe=gray!50,
    fonttitle=\bfseries, coltitle=black,
    top=4pt, bottom=4pt, left=8pt, right=8pt,
  },
}

% ─────────────────────────────────────────────────────────────────────────────
\begin{document}

% ══ PORTADA ══════════════════════════════════════════════════════════════════
\begin{center}
  {\Large\bfseries\color{azulTitulo} Universidad de Oriente/Núcleo Sucre --- Física para Bioanalistas}\\[4pt]
  {\large Informe de Evaluación Preliminar}\\[12pt]
  \rule{\linewidth}{1.2pt}\\[6pt]
  {\LARGE\bfseries krischel bravo 32449538}\\[4pt]
  \rule{\linewidth}{0.6pt}
\end{center}

% ══ FICHA TÉCNICA ════════════════════════════════════════════════════════════
\vspace{4pt}
\begin{tcolorbox}[cajaTitulo, title=Datos del informe]
\begin{tabular}{@{}llll@{}}
  \textbf{Estudiante:}  & krischel bravo 32449538  &
  \textbf{Fecha:}       & 2026-06-26 \\[4pt]
  \textbf{Archivo PDF:} & \texttt{krischel\_bravo\_32449538.pdf}  &
  \textbf{Páginas:}     & 4 \\
\end{tabular}
\end{tcolorbox}

% ══ RESULTADO GLOBAL ═════════════════════════════════════════════════════════
\vspace{6pt}
\begin{center}
  \begin{tcolorbox}[
    enhanced,
    colback=white, colframe=desempeno_bueno,
    width=0.62\linewidth,
    boxrule=2pt,
    halign=center, valign=center,
    top=8pt, bottom=8pt,
  ]
    \centering
    {\large\bfseries Puntaje sugerido}\\[6pt]
    {\Huge\bfseries\color{desempeno_bueno} 8.5/10}\\[6pt]
    {\large Nivel de desempeño: \textbf{\color{desempeno_bueno} Bueno}}
  \end{tcolorbox}
\end{center}

% ══ RESUMEN GENERAL ══════════════════════════════════════════════════════════
\section*{Resumen general}
El estudiante mostró un buen dominio de los conceptos físicos aplicados a la bioanálisis, aunque cometió algunos errores procedimentales en cálculos y manejo de unidades.

% ══ FORTALEZAS Y ÁREAS DE MEJORA ════════════════════════════════════════════
\vspace{4pt}
\begin{minipage}[t]{0.48\linewidth}
  \begin{tcolorbox}[cajaFortaleza, title={Fortalezas}]
\begin{itemize}[leftmargin=1.5em, itemsep=2pt]
  \item Correcta aplicación del principio de Arquímedes para determinar la densidad del tejido óseo.
  \item Comprensión del concepto de presión hidrostática en el contexto médico.
  \item Uso adecuado de la ecuación de Bernoulli para calcular la presión en una arteria estrechada.
  \item Capacidad para resolver problemas de flujo viscoso y resistencia en tubos.
  \item Identificación correcta de la relación entre presión y altura en fluidos.
\end{itemize}
  \end{tcolorbox}
\end{minipage}
\hfill
\begin{minipage}[t]{0.48\linewidth}
  \begin{tcolorbox}[cajaMejora, title={Areas de mejora}]
\begin{itemize}[leftmargin=1.5em, itemsep=2pt]
  \item Errores en el cálculo de la densidad del tejido óseo (confusión entre peso y masa).
  \item Falta de consistencia en el uso de unidades (ejemplo: mezcla de Pa y mmHg).
  \item Errores procedimentales en la aplicación de la ley de Poiseuille (incorrecto cálculo del caudal).
  \item Pérdida de puntos por redondeo prematuro en cálculos intermedios.
  \item Falta de claridad en la presentación de algunos pasos matemáticos.
\end{itemize}
  \end{tcolorbox}
\end{minipage}

% ══ OBSERVACIONES POR PREGUNTA ═══════════════════════════════════════════════
\section*{Observaciones por pregunta}

\begin{longtable}{>{\bfseries}p{2.8cm} p{1.6cm} p{7.0cm} p{2.8cm}}
  \toprule
  \textbf{Pregunta} & \textbf{Puntaje} & \textbf{Evaluación} & \textbf{Errores detectados} \\
  \midrule
  \endhead
  \bottomrule
  \endfoot
    \textbf{1. Presión Hidrostática y Venosa} & 2.0/2.0 & Respuesta correcta y completa. El estudiante aplicó correctamente la fórmula P = $\rho$gh y resolvió el problema con unidades adecuadas. & \textit{---} \\
    \midrule
    \textbf{2. Principio de Arquímedes} & 1.5/2.0 & La respuesta es parcialmente correcta. El estudiante aplicó el principio de Arquímedes correctamente, pero cometió un error en el cálculo de la masa del tejido óseo (confundió peso y masa). & \textit{Confundió peso y masa en el cálculo de la densidad.} \\
    \midrule
    \textbf{3. Hemodinámica (Continuidad y Bernoulli)} & 1.8/2.0 & Respuesta correcta en general. El estudiante aplicó correctamente la ecuación de Bernoulli, pero omitió la conversión de unidades en algunos pasos. & \textit{Falta de conversión de unidades en algunos pasos del cálculo.} \\
    \midrule
    \textbf{4. Ley de Poiseuille} & 1.2/2.0 & La respuesta es parcialmente correcta. El estudiante aplicó la fórmula de Poiseuille, pero cometió errores en el cálculo del radio y el caudal. & \textit{Cálculo incorrecto del radio de la arteria (usó 4 mm en lugar de 2 mm).; Error en el cálculo del caudal (usó 2 mL/s en lugar de 2,0 x 10\textasciicircum{}-7 m$^{3}$/s).} \\
    \midrule
    \textbf{5. Flujo Viscoso y Resistencia} & 2.0/2.0 & Respuesta correcta y completa. El estudiante aplicó correctamente la fórmula de resistencia y resolvió el problema con unidades adecuadas. & \textit{---} \\
    \midrule
\end{longtable}

% ══ ERRORES DETECTADOS ═══════════════════════════════════════════════════════
\section*{Análisis de errores}

\subsection*{Errores conceptuales}
\begin{itemize}[leftmargin=1.5em, itemsep=2pt]
  \item Confusión entre peso y masa en el cálculo de la densidad del tejido óseo.
  \item Uso incorrecto de unidades en algunos cálculos (ejemplo: mezcla de Pa y mmHg).
\end{itemize}

\subsection*{Errores procedimentales}
\begin{itemize}[leftmargin=1.5em, itemsep=2pt]
  \item Errores en el cálculo del radio de la arteria en la pregunta 3.
  \item Cálculo incorrecto del caudal en la pregunta 4.
  \item Pérdida de precisión en cálculos intermedios debido a redondeo prematuro.
  \item Falta de claridad en la presentación de algunos pasos matemáticos.
\end{itemize}

% ══ RECOMENDACIONES AL ESTUDIANTE ════════════════════════════════════════════
\section*{Retroalimentación para el estudiante}
\begin{tcolorbox}[cajaRecomendacion, title={Dirigido a: krischel bravo 32449538}]
El estudiante debe practicar más la conversión de unidades y el manejo de cálculos intermedios. También es importante revisar cuidadosamente los pasos matemáticos para evitar errores procedimentales. Se recomienda trabajar en la claridad y organización de las respuestas para mejorar la presentación de las soluciones.
\end{tcolorbox}

% ══ NOTA PARA EL PROFESOR ════════════════════════════════════════════════════
\section*{Nota para el profesor}
\begin{tcolorbox}[
  enhanced, breakable,
  colback=yellow!8, colframe=yellow!60!black,
  fonttitle=\bfseries, coltitle=black,
  title={Observaciones para revisión manual},
  top=4pt, bottom=4pt, left=8pt, right=8pt,
]
El examen muestra un buen nivel de comprensión conceptual, pero hay áreas donde el estudiante necesita mejorar su precisión en cálculos y manejo de unidades. Se sugiere revisar las preguntas 2 y 4 para verificar si los errores procedimentales afectaron significativamente la calificación.
\end{tcolorbox}

% ══ PIE DE INFORME ════════════════════════════════════════════════════════════
\vspace{12pt}
\begin{center}
  \footnotesize\color{gray}
  Informe generado automáticamente por el sistema de evaluación con IA (nvidia/nemotron-nano-12b-v2-vl:free) y soporte LaTex. \\
  Este documento es una evaluación \textbf{preliminar} para ser revisado por el profesor antes de ser entregado al 
  estudiante. \\
  Generado el 2026-06-26 \textbullet\ BIOANALISIS Sistema Académico de Evaluación.
  \end{center}

\end{document}
